\section{C语言：从文本文件到可执行程序的底层探秘}

\subsection{核心示例代码}
为了深入理解编译器如何将代码“分门别类”地
放入不同的内存段，我们使用以下包含各种存储
类型变量的典型示例代码进行分析：

\begin{lstlisting}[caption={内存段示例代码},label={lst:memory_segments}]
#include <stdio.h>

int globle_var1 = 10;                // 全局，已初始化
int globle_var2;                     // 全局，未初始化
static int globle_static_var1 = 20;  // 全局静态，已初始化
static int globle_static_var2;       // 全局静态，未初始化
const int globle_const_var = 30;     // 全局常量

void func(int i) {
    printf ("%d\n", i);               // 包含字符串字面量
}

int main() {
    static int static_var1 = 70;     // 局部静态，已初始化
    static int static_var2;          // 局部静态，未初始化
    const int const_var = 40;        // 局部常量
    int var1 = 30;                   // 局部变量
    int var2;                        // 局部变量
    
    // ... 运算逻辑 ...
    return 0;
}
\end{lstlisting}
\newpage


为了深入理解上述变量为何会落入不同的内存区域，
我们需要追踪这段代码经历的完整构建过程。我们称之为（\texttt{build}）。
一个完整的 C 语言构建过程包括四个主要
阶段：预处理（Preprocessing）、编
译（Compilation）、汇编（Assembly）和链接（Linking）。
下面我们逐步解析每个阶段的工作内容及其对变量存储位置的影响。

\subsubsection{预处理 (Preprocessing)}
\textbf{命令：} \texttt{gcc -E code.c -o code.i}

在这个阶段，预处理器（cpp）处理所有以 \texttt{\#} 开头的指令。对于我们的示例代码：
\begin{itemize}
    \item \textbf{头文件展开：} \texttt{\#include <stdio.h>} 被替换为标准输入输出库的完整声明内容。
    \item \textbf{注释擦除：} 所有的 \texttt{// ...} 注释被删除，留下的全是纯粹的代码。
\end{itemize}
此时，变量还只是文本形式的符号，尚未分配内存地址。

\subsubsection{编译 (Compilation)}
\textbf{命令：} \texttt{gcc -S code.i -o code.s}

编译器（cc1）将预处理后的 C 代码翻译成汇编语言。这是决定变量“命运”的关键时刻，编译器根据变量的\textbf{生命周期}（Storage Duration）、\textbf{链接属性}（Linkage）和\textbf{读写权限}（Mutability）将其放入特定的段（Segment）[cite: 58]。
\begin{itemize}
    \item \textbf{代码逻辑：} \texttt{func} 和 \texttt{main} 函数的操作被翻译成汇编指令。
    \item \textbf{数据标记：} 编译器使用汇编伪指令（Directives）标记变量所属的段。
    \item \textbf{局部变量处理：} 对于 \texttt{var1} 等局部变量，编译器不会为它们生成专门的段标签，而是生成相对于栈指针的偏移量指令。
\end{itemize}

\subsubsection{3. 汇编 (Assembly)}
\textbf{命令：} \texttt{gcc -c code.s -o code.o}

汇编器（as）将汇编代码转换为机器码，生成\textbf{可重定位目标文件}（Relocatable Object File）。此时，文件内部已经按照“段”进行了划分：

\paragraph{A. .data 段（已初始化数据段）}
[cite\_start]此段存放程序中已初始化且初值不为0的全局变量和静态变量。这些数据在程序运行前就已经确定，并存在于可执行文件中 [cite: 66, 110]。
\begin{itemize}
    \item \texttt{globle\_var1} (值: 10)
    \item \texttt{globle\_static\_var1} (值: 20)
    \item \texttt{static\_var1} (值: 70) \textbf{注意：}虽然它定义在 \texttt{main} 函数内部，但由于 \texttt{static} 关键字，它的生命周期贯穿整个程序，因此存放在数据段而非栈上。
\end{itemize}

\paragraph{B. .bss 段（未初始化数据段）}
此段存放未初始化的全局变量和静态变量。
[cite\_start]\textbf{关键特性：} .bss 段在磁盘上的目标文件中\textbf{不占用实际存储空间}，仅记录所需空间的大小。程序加载时，操作系统会将这块内存清零 [cite: 64, 110, 112]。
\begin{itemize}
    \item \texttt{globle\_var2} (默认为 0)
    \item \texttt{globle\_static\_var2} (默认为 0)
    \item \texttt{static\_var2} (默认为 0)
\end{itemize}

\paragraph{C. .rodata 段（只读数据段）}
[cite\_start]此段存放程序中的只读数据。加载到内存后，操作系统通常会将该页面标记为“只读”（Read-Only），任何写入尝试都会导致段错误（Segmentation Fault）[cite: 65, 111]。
\begin{itemize}
    \item \texttt{globle\_const\_var} (值: 30) \textendash{} 全局常量。
    \item {[cite\_start]} \texttt{"\%d\textbackslash{}n"} \textendash{} \textbf{字符串字面量}。虽然它在代码中看起来像是 \texttt{printf} 的参数，但在底层，它作为一个匿名数组存储在 .rodata 中，指令通过地址引用它 [cite: 65]。
\end{itemize}

\paragraph{D. .text 段（代码段）}
[cite\_start]此段存放编译后的机器指令（Machine Code）。该段通常是只读的，以防止程序意外修改自身的指令逻辑 [cite: 61, 105]。
\begin{itemize}
    \item \texttt{func} 函数的所有逻辑指令。
    \item \texttt{main} 函数的所有逻辑指令。
\end{itemize}

\subsubsection{4. 链接 (Linking)}
\textbf{命令：} \texttt{gcc code.o -o code}

链接器（ld）将我们的 \texttt{code.o} 与标准库合并，生成最终的\textbf{可执行文件}。
\begin{itemize}
    \item \textbf{地址重定位：} 链接器为每个段分配最终的虚拟内存地址。
    \item \textbf{符号解析：} 链接器找到 \texttt{printf} 的实现并将其地址填入调用处。
    \item \textbf{段合并：} 将多个目标文件的同类型段（如 .text）合并。
\end{itemize}

\subsection{3. 易混淆点：栈（Stack）与立即数}

并非所有变量都会被“归档”到上述的数据段中。局部变量通常在程序运行时动态创建。

\begin{itemize}
    \item \textbf{普通局部变量} (\texttt{var1}, \texttt{var2})：
    [cite\_start]它们存储在\textbf{栈（Stack）}上。在编译生成的目标文件（.o）中，找不到它们的固定地址，它们只存在于相对于栈指针（SP）的偏移量指令中 [cite: 70]。
    
    \item \textbf{局部常量} (\texttt{const\_var = 40})：
    \textbf{陷阱：} 尽管它被声明为 \texttt{const}，但因为它在函数内部，编译器通常将其视为普通的栈变量（只是编译器在编译期检查不让你修改它）。
    在底层优化中，如果数值较小（如 40），编译器甚至可能直接将其作为\textbf{立即数（Immediate Value）}嵌入到 \texttt{.text} 段的指令操作数中，根本不会在内存中分配空间。
\end{itemize}

\subsection{4. 总结对照表}

\begin{table}[h]
\centering
\begin{tabular}{|l|l|l|l|}
\hline
\textbf{变量类型} & \textbf{示例变量} & \textbf{存放位置} & \textbf{特点} \\ \hline
全局/静态 (已初始化) & \texttt{globle\_var1}, \texttt{static\_var1} & \textbf{.data} & 读写，占用磁盘空间 \\ \hline
全局/静态 (未初始化) & \texttt{globle\_var2}, \texttt{static\_var2} & \textbf{.bss} & 读写，不占磁盘空间 \\ \hline
全局常量 & \texttt{globle\_const\_var} & \textbf{.rodata} & 只读 \\ \hline
字符串字面量 & \texttt{"\%d\textbackslash{}n"} & \textbf{.rodata} & 只读，匿名存储 \\ \hline
局部变量 & \texttt{var1}, \texttt{var2} & \textbf{Stack} & 运行时分配，临时 \\ \hline
局部常量 & \texttt{const\_var} & \textbf{Stack/Text} & 视优化情况而定 \\ \hline
函数代码 & \texttt{func}, \texttt{main} & \textbf{.text} & 只读，可执行 \\ \hline
\end{tabular}
\caption{C语言变量与内存段映射关系总结}
\end{table}

\subsection{5. 为什么要进行分段？}

[cite\_start]根据《编译过程详解》，分段不仅仅是为了分类，更是为了系统效率和安全 [cite: 104]：
\begin{enumerate}
    \item {[cite\_start]} \textbf{权限管理：} 数据段可读写，代码段和只读数据段只读。这利用虚拟内存机制防止了指令被恶意或无意篡改 [cite: 105]。
    \item {[cite\_start]} \textbf{缓存效率：} 现代 CPU 的 L1 缓存通常分为指令缓存（I-Cache）和数据缓存（D-Cache）。分离存放有助于提高缓存命中率 [cite: 106]。
    \item {[cite\_start]} \textbf{内存节省：} 当多个进程运行同一个程序时，\texttt{.text} 段和 \texttt{.rodata} 段在物理内存中只需要保留一份（共享内存），只有 \texttt{.data} 和 \texttt{.bss} 需要为每个进程单独复制副本 [cite: 107]。
\end{enumerate}






































\subsection{最小示例 \texttt{HelloWorld}}%\texttt：等宽字体显示
\begin{lstlisting}[caption={最小 C 程序示例},label={lst:hello}]

  #include <stdio.h>          // 1) 预处理指令：包含头文件

int main(void) {            // 2) 主函数：程序从这里开始执行
    printf("Hello, World!\n"); // 3) 调用库函数进行输出
    return 0;               // 4) 返回状态码 0：正常结束
}
\end{lstlisting}

\subsection{\#include 与头文件}
\begin{itemize}%用于创建无序列表（项目符号列表）的环境
  \item \texttt{\#include }：\#include是预处理指令，在编译开始
  前，编译器会把被包含的文件内容“复制粘贴”进来。
  \item \texttt{stdio.h} :  \texttt{stdio.h} 中定义了常用的输入输出函数
  （如 printf, scanf）。如果事先不在预处理阶段声明定义，后续却在代码中使用了
  stdio.h 中的函数，那么编译器就会报错。
  \item 尖括号 \verb|<...>| 在系统路径查找；引号 \verb|"my.h"| 先在当前目录再到系统路径查找。
\end{itemize}

\subsection{\texttt{main} 函数（程序入口）}
\begin{lstlisting}[caption={常见main函数入口示例},label={lst:simple_main}]
int main(void)
// 或
int main(int argc, char *argv[])
\end{lstlisting}

\vspace{1em}%空出一行距离（当前字体高度）

在 C 语言中，main() 是程序的唯一入口函数，程序的执行从这里开始。

\vspace{1em}%空出一行距离（当前字体高度）


\begin{itemize}
  \item 接受：main函数的参数只能是 void 或命令行参数。究其原因是
  因为main函数是由系统启动代码调用的，系统启动时，系统只知道程序名、命令行参数和环境变量。
  
  \item 返回值：在 C 语言标准（C99、C11、C17、C23）中明确规定，main 函数必须返回一个 int 类型
的值，不能是 char、float、double、void 等其他类型，用来表示程序的执行状态。
\begin{enumerate}[label=\ding{\numexpr171+\arabic*}]% 圆圈编号，自动编号为①②③…
  \item \texttt{return 0}: 表示正常结束。
  \item \texttt{return 非 0}: 表示异常结束，具体值可自定义以表示不同错误类型。
\end{enumerate}

\end{itemize}





\subsection{语句块与分号}
\begin{itemize}
  \item 花括号 \verb|{ ... }| 构成语句块（复合语句），可包含声明与可执行语句。
  \item 大多数 C 语句以分号 \texttt{;} 结尾：表达式语句、声明（含初始化）、\texttt{return} 等。
\end{itemize}

\subsection{\texttt{printf} 与常用占位符}
\begin{lstlisting}[caption={printf 与占位符示例}]
#include <stdio.h>
int main(void) {
    int a = 42; double x = 3.14159; char c = 'A';
    printf("a=%d, x=%.2f, c=%c\n", a, x, c); // 输出 a=42, x=3.14, c=A
    return 0;
}
\end{lstlisting}




\subsection{从源代码到可执行文件：将自然语言翻译为机器语言}

C 程序从源代码到可执行文件一般经历四个主要阶段：\textbf{预处理}、\textbf{编译}、\textbf{汇编} 和 \textbf{链接}。下面以 \texttt{gcc} 为例分别说明，并给出对应命令。为便于工程化，还补充多文件与库的常见用法及常见问题。

\paragraph{（1）预处理（Preprocessing）}
主要工作：展开 \#include、宏替换（\#define）、条件编译（\#if/\#ifdef 等），并去除注释，得到纯 C 源。
\begin{verbatim}
gcc -E hello.c -o hello.i
\end{verbatim}
产物：\texttt{.i} 文件（预处理后的 C 代码）。

\paragraph{（2）编译（Compilation）}
对 \texttt{.i} 做词法/语法/语义分析与优化，生成目标架构的汇编代码。
\begin{verbatim}
gcc -S hello.i -o hello.s      % 也可直接对 .c 使用 -S
\end{verbatim}
产物：\texttt{.s} 文件（汇编源码）。

\paragraph{（3）汇编（Assembly）}
把汇编转为可重定位的目标文件，包含代码段、数据段与符号表。
\begin{verbatim}
gcc -c hello.s -o hello.o      % 也可 gcc -c hello.c -o hello.o
\end{verbatim}
产物：\texttt{.o} 文件（目标文件）。

\paragraph{（4）链接（Linking）}
将若干 \texttt{.o} 与所需库合并，进行符号解析与重定位，加入启动代码，生成可执行文件（Linux 为 ELF，Windows 为 PE）。
\begin{verbatim}
gcc hello.o -o hello
\end{verbatim}

\paragraph{（5）装载与运行（Loading \& Running）}
操作系统加载可执行文件到内存，先进入运行时入口 \texttt{\_start}，再调用用户的 \texttt{main()}；\texttt{main} 返回值作为进程退出码。

\paragraph{（6）多文件与库的基本用法}
\begin{verbatim}
gcc -c main.c -o main.o
gcc -c util.c -o util.o
gcc main.o util.o -o app              % 链接生成可执行文件

ar rcs libmylib.a foo.o bar.o         % 生成静态库
gcc main.o -L. -lmylib -o app2        % 使用静态库（库放在对象文件之后）

gcc -fPIC -c foo.c -o foo.o
gcc -shared -o libmylib.so foo.o      % 生成共享库
gcc main.o -L. -lmylib -o app3        % 运行时需能找到共享库
\end{verbatim}

\paragraph{（7）常见问题速查}
\begin{itemize}
  \item 头文件找不到：添加包含路径 \texttt{-Ipath}。
  \item 未定义引用（\texttt{undefined reference}）：缺少对应对象文件或库，或库链接顺序在前；应将 \texttt{-lxxx} 放在使用到该库的对象文件之后。
  \item 库找不到：添加库路径 \texttt{-Lpath}，并确认库名（\texttt{libxxx.a}/\texttt{libxxx.so} 对应 \texttt{-lxxx}）。
  \item 位宽/架构不匹配：确认 \texttt{-m32}/\texttt{-m64} 或交叉编译器是否正确。
\end{itemize}

\paragraph{（8）嵌入式特例}
在无操作系统的嵌入式环境中，仍有编译与链接步骤，但会使用启动文件与链接脚本生成固件映像（如 \texttt{.hex}/\texttt{.bin}），\texttt{main()} 通常为永不返回的主循环。

构建流程：
\begin{center}
\texttt{.c 源文件}
$\Rightarrow$ 预处理（展开 \#include/\#define）
$\Rightarrow$ 编译（生成 \texttt{.o/.obj}）
$\Rightarrow$ 链接（合并库与目标文件）
$\Rightarrow$ 可执行文件
\end{center}

命令行示例（以 \texttt{gcc} 为例）：
\begin{verbatim}
gcc hello.c -o hello   # 生成可执行文件
./hello                # 运行（Windows: hello.exe）
\end{verbatim}
